\section{Background}
\label{sec:background}

\paragraph{Activation steering.}
A line of recent work modifies the hidden states of a frozen language model at
inference time to push generations toward a target concept or away from an
undesired one. Inference-Time Intervention~\citep{liu2024iti} adds a learned
direction to a small set of attention heads to reduce hallucination on
TruthfulQA. Representation Engineering~\citep{zou2023repe} formulates the
intervention as a contrastive readout-and-edit problem on a chosen layer's
residual stream and shows the signal generalizes across many concepts.
Subsequent work refines where and how to read the contrastive signal: difference
of means across paired prompt sets~\citep{li2024fewshot}, projected gradient
descent across multiple residual layers~\citep{turner2023acteng}, and
attention-guided readouts that weight token positions by attention rather than
by fixed slots~\citep{tan2024attentionguided}.

All of these methods share an architectural assumption that the per-token
hidden-state direction is the natural intervention target.

\paragraph{Sparse MoE language models.}
\mixtral~\citep{jiang2024mixtral} and \olmoe~\citep{muennighoff2024olmoe} are
sparse mixture-of-experts language models in which each MoE layer routes every
token to a small subset (typically two) of $E$ feed-forward experts using a
learned router. Sparse MoEs decouple parameter count from per-token compute
and have become a common architectural choice for open-weight models trained
beyond a few billion active parameters.

\paragraph{\steermoe.}
\steermoe~\citep{wang2024steermoe} argues that the natural intervention site in
a sparse MoE is not a hidden-state direction but the router itself. Given a
contrastive prompt pair (one prompt that elicits the target behavior, one that
does not), \steermoe collects per-token router logits across MoE layers,
records which experts were selected by top-$k$ routing, computes a
risk-difference table over $(\text{layer}, \text{expert})$ cells, picks the
strongest positive and negative experts, and adds a fixed bias to those
experts' router logits at generation time. The empirical result on the original
small instruct MoEs is that this routing-bias intervention can move generations
toward the target concept in a way that hidden-state steering on the same
checkpoint cannot.

\paragraph{What this paper asks.}
\steermoe's published results were obtained on a particular family of small
instruct MoEs. Our question is whether the recipe --- the readout (matched
target span, risk-difference table, top-$k$ expert selection) and the
intervention (additive router-logit bias on selected experts) --- transfers to
\mixtral on a different concept-transfer dataset. We adopt an explicitly
diagnostic stance: we are not trying to set state of the art on a benchmark,
and we do not claim a method contribution. We are checking whether a
faithful re-implementation produces the kind of concept-specific shifts the
upstream paper describes when held to a question-only test in which the
concept prefix is omitted at evaluation time.
